%2multibyte Version: 5.50.0.2960 CodePage: 65001
% !TeX root = ./Slides_Lez_13_Python_Asset_Allocation_ALM.tex
% ------------------------------------------------------------
% Macro matematiche ufficiali del progetto MQF
% Coerenti con MQF_Notazione.tex
% ------------------------------------------------------------
% Fallback utile per compilazione esterna se qualche titolo resta scritto come \QTR.
% Scientific WorkPlace usa \QTR internamente; in Beamer \QTR{frametitle}{Titolo}
% deve espandersi in \frametitle{Titolo}.

\documentclass[handout]{beamer}%
\usepackage{mathpazo}
%\usepackage{hyperref}
\usepackage{multimedia}
\usepackage{amsmath}
\usepackage{amsfonts}
\usepackage{amssymb}
\usepackage{graphicx}%
\setcounter{MaxMatrixCols}{30}
%TCIDATA{OutputFilter=latex2.dll}
%TCIDATA{Version=5.50.0.2960}
%TCIDATA{Codepage=65001}
%TCIDATA{CSTFile=beamer.cst}
%TCIDATA{Created=Saturday, May 16, 2026 09:09:59}
%TCIDATA{LastRevised=Saturday, May 16, 2026 11:24:58}
%TCIDATA{<META NAME="GraphicsSave" CONTENT="32">}
%TCIDATA{<META NAME="SaveForMode" CONTENT="2">}
%TCIDATA{BibliographyScheme=Manual}
%TCIDATA{<META NAME="DocumentShell" CONTENT="Other Documents\SW\Slides - Beamer">}
%BeginMSIPreambleData
\graphicspath{{../graphics/}}
\providecommand{\U}[1]{\protect\rule{.1in}{.1in}}
%EndMSIPreambleData
\newenvironment{stepenumerate}{\begin{enumerate}[<+->]}{\end{enumerate}}
\newenvironment{stepitemize}{\begin{itemize}[<+->]}{\end{itemize} }
\newenvironment{stepenumeratewithalert}{\begin{enumerate}[<+-| alert@+>]}{\end{enumerate}}
\newenvironment{stepitemizewithalert}{\begin{itemize}[<+-| alert@+>]}{\end{itemize} }
\usetheme{Madrid}


\setbeamertemplate{footline}
{
  \leavevmode%
  \hbox{%
    \begin{beamercolorbox}[wd=.20\paperwidth,ht=2.25ex,dp=1ex,center]
      {author in head/foot}%
      \usebeamerfont{author in head/foot}\insertshortauthor
    \end{beamercolorbox}%
    \begin{beamercolorbox}[wd=.60\paperwidth,ht=2.25ex,dp=1ex,center]
      {title in head/foot}%
      \usebeamerfont{title in head/foot}\insertshorttitle
    \end{beamercolorbox}%
    \begin{beamercolorbox}[wd=.20\paperwidth,ht=2.25ex,dp=1ex,right]
      {date in head/foot}%
      \usebeamerfont{date in head/foot}%
      \insertshortdate{}\hspace*{2em}%
      \insertframenumber{} / \inserttotalframenumber\hspace*{2ex}
    \end{beamercolorbox}%
  }%
  \vskip0pt%
}




\providecommand{\R}{\mathbb{R}}
\providecommand{\N}{\mathbb{N}}
\providecommand{\E}{\mathbb{E}}
\providecommand{\Prob}{\mathbb{P}}
\providecommand{\Q}{\mathbb{Q}}
\providecommand{\F}{\mathcal{F}}
\providecommand{\B}{\mathcal{B}}
\providecommand{\G}{\mathcal{G}}
\providecommand{\1}{\mathbf{1}}
\providecommand{\Var}{\operatorname{Var}}
\providecommand{\Cov}{\operatorname{Cov}}
\providecommand{\VaR}{\operatorname{VaR}}
\providecommand{\CVaR}{\operatorname{CVaR}}
\providecommand{\argmin}{\operatorname*{arg\,min}}
\providecommand{\argmax}{\operatorname*{arg\,max}}
\providecommand{\EVPI}{\operatorname{EVPI}}
\providecommand{\VSS}{\operatorname{VSS}}
\providecommand{\QTR}[2]{\csname #1\endcsname{#2}}
%BeginMSIPreambleData
\ifx\pdfoutput\relax\let\pdfoutput=\undefined\fi
\newcount\msipdfoutput
\ifx\pdfoutput\undefined\else
\ifcase\pdfoutput\else
\msipdfoutput=1
\ifx\paperwidth\undefined\else
\ifdim\paperheight=0pt\relax\else\pdfpageheight\paperheight\fi
\ifdim\paperwidth=0pt\relax\else\pdfpagewidth\paperwidth\fi
\fi\fi\fi
%EndMSIPreambleData

%******************
% Slide 1
%******************

\title[Lez. 13 - Python, programmazione lineare e ALM]
{Lezione 13 - Applicazione in Python: programmazione lineare e ALM deterministico}

\subtitle{Metodi Quantitativi per la Finanza}

\author[P. Falbo, C. Pelizzari]
{{\LARGE Paolo Falbo, Cristian Pelizzari}}

\institute[]{{\large Dipartimento di Economia e Management}}

\date[]{}

\titlegraphic{
\includegraphics[width=3.3cm]{../graphics/logo-Unibs-esteso.png}
}

\begin{document}

\maketitle


\section{Apertura e obiettivi}

\subsection{Dalla teoria alla soluzione computazionale}

%******************
% Slide 2
%******************

%TCIMACRO{\TeXButton{BeginFrame}{\begin{frame}}}%
%BeginExpansion
\begin{frame}%
%EndExpansion

\QTR{frametitle}{Dalla teoria alla soluzione computazionale}

%TCIMACRO{\TeXButton{Transparency}{\setbeamercovered{transparent=20}}}%
%BeginExpansion
\setbeamercovered{transparent=20}%
%EndExpansion

La Lezione 13 consolida operativamente i modelli introdotti nelle due lezioni
precedenti.

\vspace{0.4cm}

\[
\text{Lezione 11}
\quad\longrightarrow\quad
\text{Lezione 12}
\quad\longrightarrow\quad
\text{Lezione 13}
\]

\vspace{-0.2cm}

\[
\text{progr. lin.}
\rightarrow
\text{ALM deterministico}
\rightarrow
\text{implementazione e analisi computaz.}
\]

\vspace{0.4cm}

\begin{stepitemize}

\item La lezione 11 ha introdotto la scelta ottima sotto vincoli, la dualit\`a
e i valori marginali.

\item La lezione 12 ha introdotto la dimensione temporale dei cash flow,
i bilanci di liquidit\`a e il funding.

\item In questa lezione, il modello viene tradotto nella forma richiesta dal
solver, risolto, controllato e sottoposto ad analisi what-if.

\end{stepitemize}

\vspace{0.2cm}

\textbf{Idea guida:} non introduciamo una nuova classe di modelli; rendiamo
operativo il modello gi\`a costruito.

%TCIMACRO{\TeXButton{Transition: Box Out}{\transboxout}}%
%BeginExpansion
\transboxout%
%EndExpansion
%TCIMACRO{\TeXButton{EndFrame}{\end{frame}}}%
%BeginExpansion
\end{frame}%
%EndExpansion

\subsection{Obiettivi e struttura della lezione}

%******************
% Slide 3
%******************

%TCIMACRO{\TeXButton{BeginFrame}{\begin{frame}}}%
%BeginExpansion
\begin{frame}%
%EndExpansion

\QTR{frametitle}{Obiettivi e struttura della lezione}

%TCIMACRO{\TeXButton{Transparency}{\setbeamercovered{transparent=20}}}%
%BeginExpansion
\setbeamercovered{transparent=20}%
%EndExpansion

La lezione persegue tre obiettivi:

\begin{stepenumerate}

\item [1.] \textbf{Un piccolo approfondimento Python}

Isolare una struttura del codice utile per costruire e analizzare
ripetutamente un modello parametrico.

\item [2.] \textbf{Il Caso Aula}

Tradurre un problema di ALM deterministico nella forma richiesta dal solver,
verificare la soluzione e impostare benchmark, sensitività, frontiera di
fattibilit\`a e scenari.

\item [3.] \textbf{Il Caso TakeHome}

Presentare un nuovo problema finanziario al quale trasferire autonomamente
lo stesso metodo di lavoro, senza svilupparne la soluzione in aula.

\end{stepenumerate}

%\vspace{0.3cm}

La lezione sar\`a articolata in due momenti:

\vspace{-0.5cm}

\[
\text{inquadramento con le slide}
\quad\longrightarrow\quad
\text{sviluppo del Caso Aula in VS Code}.
\]

%\vspace{0.2cm}

\textbf{Priorit\`a:} il Caso Aula costituisce il nucleo principale della
lezione; le slide ne preparano lo sviluppo computazionale senza sostituirlo.

%TCIMACRO{\TeXButton{Transition: Box Out}{\transboxout}}%
%BeginExpansion
\transboxout%
%EndExpansion
%TCIMACRO{\TeXButton{EndFrame}{\end{frame}}}%
%BeginExpansion
\end{frame}%
%EndExpansion

\section{Un piccolo focus Python}

\subsection{Dal codice ripetuto a una funzione}

%******************
% Slide 4
%******************

%TCIMACRO{\TeXButton{BeginFrame}{\begin{frame}}}%
%BeginExpansion
\begin{frame}%
%EndExpansion

\QTR{frametitle}{Dal codice ripetuto a una funzione}

%TCIMACRO{\TeXButton{Transparency}{\setbeamercovered{transparent=20}}}%
%BeginExpansion
\setbeamercovered{transparent=20}%
%EndExpansion

Nel Caso Aula, lo stesso modello deve essere risolto pi\`u volte al variare
di alcuni parametri.

\vspace{0.3cm}

Una possibilit\`a sarebbe ripetere ogni volta le stesse istruzioni.
La soluzione pi\`u naturale \`e invece racchiudere il procedimento in una
\textbf{funzione}.

%\vspace{0.4cm}

\[
(\theta,\phi,\psi)
\quad\longrightarrow\quad
\text{costruzione del problema}
\quad\longrightarrow\quad
\text{soluzione}.
\]

\vspace{0.4cm}

In Python, la struttura essenziale \`e:

%\vspace{0.2cm}

\begin{center}
\texttt{def solve\_model(theta, phi, psi):}

%\vspace{0.15cm}

\hspace{0.8cm}\texttt{...}

%\vspace{0.15cm}

\hspace{0.8cm}\texttt{return result}
\end{center}

%\vspace{0.3cm}

\textbf{Idea chiave:} la funzione rende il modello ripetibile senza
duplicare la logica di costruzione e soluzione.

%TCIMACRO{\TeXButton{Transition: Box Out}{\transboxout}}%
%BeginExpansion
\transboxout%
%EndExpansion
%TCIMACRO{\TeXButton{EndFrame}{\end{frame}}}%
%BeginExpansion
\end{frame}%
%EndExpansion


\subsection{Input, elaborazione e output}

%******************
% Slide 5
%******************

%TCIMACRO{\TeXButton{BeginFrame}{\begin{frame}}}%
%BeginExpansion
\begin{frame}%
%EndExpansion

\QTR{frametitle}{Input, elaborazione e output}

%TCIMACRO{\TeXButton{Transparency}{\setbeamercovered{transparent=20}}}%
%BeginExpansion
\setbeamercovered{transparent=20}%
%EndExpansion

Una funzione pu\`o essere letta come una trasformazione organizzata tra
\textbf{input} e \textbf{output}.

%\vspace{0.4cm}

\[
\boxed{\theta,\phi,\psi}
\quad\longrightarrow\quad
\boxed{\text{elaborazione}}
\quad\longrightarrow\quad
\boxed{\text{result}}
\]

%\vspace{0.4cm}

Nel nostro caso, l'elaborazione comprende:

\begin{stepenumerate}

\item [1.] costruire i coefficienti che dipendono dai parametri;

\item [2.] costruire uguaglianze, disuguaglianze e bound;

\item [3.] chiamare il solver;

\item [4.] restituire il risultato della soluzione.

\end{stepenumerate}

%\vspace{0.3cm}

L'oggetto restituito dal solver contiene poi le informazioni necessarie per
analizzare:

\vspace{-0.3cm}

\[
\text{soluzione}
\;,\qquad
\text{valore obiettivo}
\;,\qquad
\text{fattibilit\`a}
\;,\qquad
\text{slack e marginali}.
\]

%\vspace{0.2cm}

\textbf{Collegamento metodologico:} la struttura della funzione riflette la
scomposizione del problema in tappe input--output.

%TCIMACRO{\TeXButton{Transition: Box Out}{\transboxout}}%
%BeginExpansion
\transboxout%
%EndExpansion
%TCIMACRO{\TeXButton{EndFrame}{\end{frame}}}%
%BeginExpansion
\end{frame}%
%EndExpansion

\subsection{Da una soluzione a una famiglia di soluzioni}

%******************
% Slide 6
%******************

%TCIMACRO{\TeXButton{BeginFrame}{\begin{frame}}}%
%BeginExpansion
\begin{frame}%
%EndExpansion

\QTR{frametitle}{Da una soluzione a una famiglia di soluzioni}

%TCIMACRO{\TeXButton{Transparency}{\setbeamercovered{transparent=20}}}%
%BeginExpansion
\setbeamercovered{transparent=20}%
%EndExpansion

La funzione parametrica permette di utilizzare lo stesso modello per
domande quantitative diverse.

%\vspace{0.4cm}

\begin{stepenumerate}

\item [1.] \textbf{Benchmark}
\vspace{-0.2cm}
\[
\texttt{solve\_model(1, 1, 1)}
\]

fornisce la soluzione del caso di riferimento.

\item [2.] \textbf{Sensitivit\`a univariata}
\vspace{-0.2cm}
\[
\texttt{solve\_model(theta, 1, 1)}
\]

permette di studiare come cambia la soluzione al variare di un solo driver.

\item [3.] \textbf{Scenari congiunti}
\vspace{-0.2cm}
\[
\texttt{solve\_model(theta, phi, psi)}
\]

permette di modificare simultaneamente pi\`u driver.

\end{stepenumerate}

\vspace{-0.2cm}

\[
\boxed{\text{una funzione}}
\rightarrow
\boxed{\text{una famiglia di problemi LP}}
\rightarrow
\boxed{\text{una famiglia di soluzioni}}
\]

%\vspace{0.2cm}

\textbf{Passaggio successivo:} identifichiamo ora il modello ALM che la
funzione dovr\`a costruire e risolvere.

%TCIMACRO{\TeXButton{Transition: Box Out}{\transboxout}}%
%BeginExpansion
\transboxout%
%EndExpansion
%TCIMACRO{\TeXButton{EndFrame}{\end{frame}}}%
%BeginExpansion
\end{frame}%
%EndExpansion

\section{Caso Aula: modello e preparazione}

\subsection{Il problema ALM}

%******************
% Slide 7
%******************

%TCIMACRO{\TeXButton{BeginFrame}{\begin{frame}}}%
%BeginExpansion
\begin{frame}%
%EndExpansion

\QTR{frametitle}{Caso Aula: Silicon Valley Bank 2023}

%TCIMACRO{\TeXButton{Transparency}{\setbeamercovered{transparent=20}}}%
%BeginExpansion
\setbeamercovered{transparent=20}%
%EndExpansion

Consideriamo un modello stilizzato di Asset--Liability Management ispirato
al caso Silicon Valley Bank.

%\vspace{0.3cm}

La banca deve scegliere un piano di allocazione degli attivi e di gestione
della liquidit\`a su quattro date, potendo ricorrere anche a funding esterno.

%\vspace{0.4cm}

La domanda quantitativa centrale \`e:

%\vspace{0.2cm}

\begin{center}
\textbf{come cambia il piano ALM ottimo quando cambiano
fabbisogni di liquidit\`a, capacit\`a di funding e costo del funding?}
\end{center}

%\vspace{0.4cm}

I tre driver dell'analisi sono:

\[
\begin{array}{ccl}
\theta &:& \text{fabbisogni di liquidit\`a},\\[0.1cm]
\phi   &:& \text{capacit\`a di funding},\\[0.1cm]
\psi   &:& \text{costo del funding}.
\end{array}
\]

\vspace{0.2cm}

Una seconda domanda riguarda la \textbf{fattibilit\`a}: fino a quale livello
possono aumentare i fabbisogni iniziali prima che non esista pi\`u un piano
ammissibile?

%TCIMACRO{\TeXButton{Transition: Box Out}{\transboxout}}%
%BeginExpansion
\transboxout%
%EndExpansion
%TCIMACRO{\TeXButton{EndFrame}{\end{frame}}}%
%BeginExpansion
\end{frame}%
%EndExpansion

\subsection{Dal modello statico all'ALM temporale}

%******************
% Slide 8
%******************

%TCIMACRO{\TeXButton{BeginFrame}{\begin{frame}}}%
%BeginExpansion
\begin{frame}%
%EndExpansion

\QTR{frametitle}{Dalla Lezione 11 alla Lezione 12}

%TCIMACRO{\TeXButton{Transparency}{\setbeamercovered{transparent=20}}}%
%BeginExpansion
\setbeamercovered{transparent=20}%
%EndExpansion

Il Caso Aula combina due livelli gi\`a sviluppati nel corso.

%\vspace{0.4cm}

\begin{columns}[T]

\begin{column}{0.46\textwidth}

\begin{center}
\textbf{Lezione 11}
\end{center}

\[
x=(x_1,\ldots,x_4)'
\]

\begin{itemize}

\item allocazione degli attivi;

\item funzione obiettivo;

\item vincoli;

\item soluzione ottima;

\item dualit\`a e valori marginali.

\end{itemize}

\end{column}

\begin{column}{0.46\textwidth}

\begin{center}
\textbf{Lezione 12}
\end{center}

\[
(x,g,u)
\]

\begin{itemize}

\item cash flow distribuiti nel tempo;

\item bilanci di liquidit\`a;

\item liquidit\`a riportata;

\item funding esterno;

\item valori marginali temporali.

\end{itemize}

\end{column}

\end{columns}

%\vspace{0.3cm}

\[
\boxed{\text{La classe matematica resta la programmazione lineare}}
\]

ma la struttura economica diventa multiperiodale.

%TCIMACRO{\TeXButton{Transition: Box Out}{\transboxout}}%
%BeginExpansion
\transboxout%
%EndExpansion
%TCIMACRO{\TeXButton{EndFrame}{\end{frame}}}%
%BeginExpansion
\end{frame}%
%EndExpansion

\subsection{La dimensione temporale}

%******************
% Slide 9
%******************

%TCIMACRO{\TeXButton{BeginFrame}{\begin{frame}}}%
%BeginExpansion
\begin{frame}%
%EndExpansion

\QTR{frametitle}{Il tempo entra nel modello}

%TCIMACRO{\TeXButton{Transparency}{\setbeamercovered{transparent=20}}}%
%BeginExpansion
\setbeamercovered{transparent=20}%
%EndExpansion

Nel modello ALM, la liquidit\`a deve essere disponibile alla data in cui nasce
il fabbisogno.

%\vspace{0.5cm}

\[
\begin{array}{ccccccc}
t=1 & & t=2 & & t=3 & & t=4\\
\bullet & \longrightarrow & \bullet & \longrightarrow &
\bullet & \longrightarrow & \bullet
\end{array}
\]

%\vspace{0.4cm}

A ogni data, intervengono:

\begin{stepitemize}

\item i cash flow generati dagli attivi, \(a_t'x\);

\item il fabbisogno di liquidit\`a, \(d_t\);

\item l'eventuale funding esterno, \(u_t\);

\item la liquidit\`a residua trasferita alla data successiva, \(g_t\).

\end{stepitemize}

%\vspace{0.4cm}

\[
\boxed{
\text{dal ``quanto'' al ``quanto e quando''}
}
\]

%\vspace{0.2cm}

Il problema non riguarda quindi soltanto \textbf{quanto} cash flow viene
generato, ma anche \textbf{quando} viene generato.

%TCIMACRO{\TeXButton{Transition: Box Out}{\transboxout}}%
%BeginExpansion
\transboxout%
%EndExpansion
%TCIMACRO{\TeXButton{EndFrame}{\end{frame}}}%
%BeginExpansion
\end{frame}%
%EndExpansion



\subsection{Gli oggetti del modello}

%******************
% Slide 10
%******************

%TCIMACRO{\TeXButton{BeginFrame}{\begin{frame}}}%
%BeginExpansion
\begin{frame}%
%EndExpansion

\QTR{frametitle}{Le variabili decisionali}

%TCIMACRO{\TeXButton{Transparency}{\setbeamercovered{transparent=20}}}%
%BeginExpansion
\setbeamercovered{transparent=20}%
%EndExpansion

Il piano ALM comprende tre gruppi di variabili decisionali.

\vspace{0.4cm}

\begin{description}

\item[\textbf{Allocazione degli attivi}]
\[
x=(x_1,x_2,x_3,x_4)'.
\]

\(x_j\) indica l'ammontare investito nella classe di attivo \(j\).

\vspace{0.2cm}

\item[\textbf{Liquidit\`a riportata}]
\[
g=(g_1,g_2,g_3,g_4)'.
\]

\(g_t\) rappresenta la liquidit\`a residua dopo il fabbisogno alla data \(t\).

\vspace{0.2cm}

\item[\textbf{Funding esterno}]
\[
u=(u_1,u_2,u_3)'.
\]

\(u_t\) rappresenta il funding ottenuto alla data \(t\), disponibile
immediatamente e rimborsato alla data successiva.

\end{description}

\vspace{0.3cm}

Il modello contiene complessivamente $4+4+3=\boxed{11}$ variabili decisionali.

%TCIMACRO{\TeXButton{Transition: Box Out}{\transboxout}}%
%BeginExpansion
\transboxout%
%EndExpansion
%TCIMACRO{\TeXButton{EndFrame}{\end{frame}}}%
%BeginExpansion
\end{frame}%
%EndExpansion

\subsection{Funzione obiettivo e vincoli}

%******************
% Slide 11
%******************

%TCIMACRO{\TeXButton{BeginFrame}{\begin{frame}}}%
%BeginExpansion
\begin{frame}%
%EndExpansion

\QTR{frametitle}{La funzione obiettivo}

%TCIMACRO{\TeXButton{Transparency}{\setbeamercovered{transparent=20}}}%
%BeginExpansion
\setbeamercovered{transparent=20}%
%EndExpansion

Il piano viene scelto massimizzando il valore economico complessivo:

\[
\max_{x,g,u}
\;
\Pi(x,g,u;\psi)
=
c'x
+
\sum_{t=1}^{3} r_t^g g_t
-
\sum_{t=1}^{3}\psi r_t^u u_t.
\]

%\vspace{0.4cm}

La funzione obiettivo combina tre componenti:

\begin{stepenumerate}

\item [1.] \(c'x\): rendimento economico dell'allocazione degli attivi;

\item [2.] \(\displaystyle\sum_{t=1}^{3} r_t^g g_t\):
remunerazione della liquidit\`a riportata;

\item [3.] \(\displaystyle-\sum_{t=1}^{3}\psi r_t^u u_t\):
costo del funding esterno.

\end{stepenumerate}

%\vspace{0.3cm}

\[
\boxed{\psi}
\]

modifica quindi direttamente il costo economico del funding.

%TCIMACRO{\TeXButton{Transition: Box Out}{\transboxout}}%
%BeginExpansion
\transboxout%
%EndExpansion
%TCIMACRO{\TeXButton{EndFrame}{\end{frame}}}%
%BeginExpansion
\end{frame}%
%EndExpansion

%******************
% Slide 12
%******************

%TCIMACRO{\TeXButton{BeginFrame}{\begin{frame}}}%
%BeginExpansion
\begin{frame}%
%EndExpansion

\QTR{frametitle}{Bilanci di liquidit\`a: la struttura ricorsiva}

%TCIMACRO{\TeXButton{Transparency}{\setbeamercovered{transparent=20}}}%
%BeginExpansion
\setbeamercovered{transparent=20}%
%EndExpansion

Alla prima data:
\vspace{-0.1cm}
\[
a_1'x+u_1-g_1=d_1(\theta).
\]

%\vspace{0.3cm}

Alle date successive, compare una struttura ricorsiva:
\vspace{-0.1cm}
\[
a_t'x
+
(1+r_{t-1}^g)g_{t-1}
+
u_t
-
(1+\psi r_{t-1}^u)u_{t-1}
-
g_t
=
d_t(\theta),
\qquad t=2,3.
\]

\vspace{0.2cm}

La logica \`e sempre la stessa:

$$
\text{cash flow}
+
\text{liquidit\`a riportata}
+
\text{nuovo funding}
$$
\vspace{-0.6cm}
$$
=
$$
\vspace{-0.6cm}
$$
\text{fabbisogno}
+
\text{rimborso}
+
\text{liquidit\`a residua}
$$

\vspace{0.3cm}

\textbf{Punto chiave:} il funding ottenuto alla data \(t\) migliora la
liquidit\`a corrente, ma genera un rimborso alla data successiva.

%TCIMACRO{\TeXButton{Transition: Box Out}{\transboxout}}%
%BeginExpansion
\transboxout%
%EndExpansion
%TCIMACRO{\TeXButton{EndFrame}{\end{frame}}}%
%BeginExpansion
\end{frame}%
%EndExpansion

%******************
% Slide 13
%******************

%TCIMACRO{\TeXButton{BeginFrame}{\begin{frame}}}%
%BeginExpansion
\begin{frame}%
%EndExpansion

\QTR{frametitle}{Gli altri vincoli}

%TCIMACRO{\TeXButton{Transparency}{\setbeamercovered{transparent=20}}}%
%BeginExpansion
\setbeamercovered{transparent=20}%
%EndExpansion

Ai bilanci temporali, si aggiungono gli altri vincoli del piano ALM.

%\vspace{0.3cm}

\begin{stepenumerate}

\item [1.] \textbf{Risorse iniziali}
\vspace{-0.1cm}
\[
\sum_{j=1}^{4}x_j=A_0=100.
\]

\item [2.] \textbf{Limite prudenziale}
\vspace{-0.1cm}
\[
\ell'x\le K,
\qquad K=7.
\]

\item [3.] \textbf{Concentrazione}
\vspace{-0.1cm}
\[
x_3\le30.
\]

\item [4.] \textbf{Capacit\`a di funding}
\vspace{-0.1cm}
\[
0\le u_t\le\phi\bar u_t.
\]

\item [5.] \textbf{Non negativit\`a}
\vspace{-0.1cm}
\[
x_j\ge0,
\qquad
g_t\ge0.
\]

\end{stepenumerate}

%\vspace{0.2cm}

Tutte le relazioni sono lineari: il problema resta quindi un problema di
\textbf{programmazione lineare}.

%TCIMACRO{\TeXButton{Transition: Box Out}{\transboxout}}%
%BeginExpansion
\transboxout%
%EndExpansion
%TCIMACRO{\TeXButton{EndFrame}{\end{frame}}}%
%BeginExpansion
\end{frame}%
%EndExpansion

\subsection{Dalla formulazione matematica al solver}

%******************
% Slide 14
%******************

%TCIMACRO{\TeXButton{BeginFrame}{\begin{frame}}}%
%BeginExpansion
\begin{frame}%
%EndExpansion

\QTR{frametitle}{Il vettore del solver}

%TCIMACRO{\TeXButton{Transparency}{\setbeamercovered{transparent=20}}}%
%BeginExpansion
\setbeamercovered{transparent=20}%
%EndExpansion

Per costruire il problema numerico, dobbiamo raccogliere tutte le variabili
decisionali in un unico vettore.

\vspace{0.3cm}

Indichiamo localmente con \(y\) il vettore utilizzato nella forma solver:
\vspace{-0.1cm}
\[
y=
(x_1,x_2,x_3,x_4,
g_1,g_2,g_3,g_4,
u_1,u_2,u_3)'.
\]

%\vspace{0.4cm}

% La struttura per blocchi \`e:
% \vspace{-0.1cm}
% \[
% \underbrace{x_1,\ldots,x_4}_{\text{attivi}}
% \qquad
% \underbrace{g_1,\ldots,g_4}_{\text{liquidit\`a}}
% \qquad
% \underbrace{u_1,\ldots,u_3}_{\text{funding}}.
% \]

% \vspace{0.4cm}

Il significato economico delle variabili non cambia, ma il solver le tratta
come componenti di un unico vettore numerico.

\vspace{0.2cm}

\[
\boxed{\text{modello economico}}
\quad\longrightarrow\quad
\boxed{\text{vettore decisionale numerico}}
\]

%TCIMACRO{\TeXButton{Transition: Box Out}{\transboxout}}%
%BeginExpansion
\transboxout%
%EndExpansion
%TCIMACRO{\TeXButton{EndFrame}{\end{frame}}}%
%BeginExpansion
\end{frame}%
%EndExpansion

%******************
% Slide 15
%******************

%TCIMACRO{\TeXButton{BeginFrame}{\begin{frame}}}%
%BeginExpansion
\begin{frame}%
%EndExpansion

\QTR{frametitle}{Dalla formulazione agli input del solver}

%TCIMACRO{\TeXButton{Transparency}{\setbeamercovered{transparent=20}}}%
%BeginExpansion
\setbeamercovered{transparent=20}%
%EndExpansion

Il solver richiede quattro oggetti fondamentali.

\vspace{-0.1cm}

\[
\begin{array}{ccl}
\text{funzione obiettivo}
&\longrightarrow&
\texttt{c}
\\[0.25cm]
\text{uguaglianze}
&\longrightarrow&
\texttt{A\_eq},\ \texttt{b\_eq}
\\[0.25cm]
\text{disuguaglianze}
&\longrightarrow&
\texttt{A\_ub},\ \texttt{b\_ub}
\\[0.25cm]
\text{limiti sulle variabili}
&\longrightarrow&
\texttt{bound}
\end{array}
\]

%\vspace{-0.1cm}

La forma risolta da \texttt{linprog} \`e:
\vspace{-0.1cm}

\[
\min_y f'y
\]

soggetto a
\vspace{-0.1cm}

\[
A_{ub}y\le b_{ub},
\qquad
A_{eq}y=b_{eq}.
\]

\vspace{-0.1cm}

\textbf{Idea chiave:} la programmazione consiste nel tradurre correttamente
gli oggetti del modello in questi input, non nel modificare il modello.

%TCIMACRO{\TeXButton{Transition: Box Out}{\transboxout}}%
%BeginExpansion
\transboxout%
%EndExpansion
%TCIMACRO{\TeXButton{EndFrame}{\end{frame}}}%
%BeginExpansion
\end{frame}%
%EndExpansion

%******************
% Slide 16
%******************

%TCIMACRO{\TeXButton{BeginFrame}{\begin{frame}}}%
%BeginExpansion
\begin{frame}%
%EndExpansion

\QTR{frametitle}{Una sola inversione di segno}

%TCIMACRO{\TeXButton{Transparency}{\setbeamercovered{transparent=20}}}%
%BeginExpansion
\setbeamercovered{transparent=20}%
%EndExpansion

Il modello economico massimizza:
\vspace{-0.1cm}
\[
\max_y p'y.
\]

\vspace{0.4cm}

Poich\'e \texttt{linprog} minimizza, definiamo:
\vspace{-0.1cm}

\[
f=-p.
\]

Il solver risolve quindi:
\vspace{-0.1cm}

\[
\min_y f'y.
\]

%\vspace{0.4cm}

Una volta ottenuta la soluzione \(y^*\):
\vspace{-0.1cm}

\[
\Pi^*=p'y^*=-f'y^*.
\]

%\vspace{0.4cm}

\[
\boxed{
\text{il cambio di segno si effettua una sola volta}
}
\]

%TCIMACRO{\TeXButton{Transition: Box Out}{\transboxout}}%
%BeginExpansion
\transboxout%
%EndExpansion
%TCIMACRO{\TeXButton{EndFrame}{\end{frame}}}%
%BeginExpansion
\end{frame}%
%EndExpansion

%******************
% Slide 17
%******************

%TCIMACRO{\TeXButton{BeginFrame}{\begin{frame}}}%
%BeginExpansion
\begin{frame}%
%EndExpansion

\QTR{frametitle}{Tre parametri, tre punti di ingresso}

%TCIMACRO{\TeXButton{Transparency}{\setbeamercovered{transparent=20}}}%
%BeginExpansion
\setbeamercovered{transparent=20}%
%EndExpansion

I tre driver modificano parti diverse del problema.

%\vspace{0.5cm}

\[
\begin{array}{ccl}
\theta
&\longrightarrow&
b_{eq}
\\[0.3cm]
\phi
&\longrightarrow&
\texttt{bound}
\\[0.3cm]
\psi
&\longrightarrow&
f\ \text{e}\ A_{eq}
\end{array}
\]

%\vspace{0.5cm}

In particolare:

\[
d(\theta)=(30\theta,30\theta,22,15)',
\]
\vspace{-0.2cm}
\[
\bar u_t(\phi)=\phi\bar u_t.
\]

%\vspace{0.3cm}

Per \(\psi\), occorre ricordare due effetti:

\begin{itemize}

\item modifica il costo del funding nella funzione obiettivo;

\item modifica il rimborso del funding nei bilanci temporali.

\end{itemize}

\vspace{0.2cm}

\textbf{Punto chiave:} parametrizzare il modello significa sapere
esattamente quale oggetto del solver deve cambiare.

%TCIMACRO{\TeXButton{Transition: Box Out}{\transboxout}}%
%BeginExpansion
\transboxout%
%EndExpansion
%TCIMACRO{\TeXButton{EndFrame}{\end{frame}}}%
%BeginExpansion
\end{frame}%
%EndExpansion

%******************
% Slide 18
%******************

%TCIMACRO{\TeXButton{BeginFrame}{\begin{frame}}}%
%BeginExpansion
\begin{frame}%
%EndExpansion

\QTR{frametitle}{Che cosa faremo in VS Code}

%TCIMACRO{\TeXButton{Transparency}{\setbeamercovered{transparent=20}}}%
%BeginExpansion
\setbeamercovered{transparent=20}%
%EndExpansion

A questo punto, il modello \`e pronto per essere costruito e analizzato
numericamente.

%\vspace{0.4cm}

\[
\boxed{\text{benchmark}}
\longrightarrow
\boxed{\text{controlli}}
\longrightarrow
\boxed{\text{sensitivit\`a}}
\longrightarrow
\boxed{\text{frontiera}}
\longrightarrow
\boxed{\text{scenari}}
\]

%\vspace{0.5cm}

Nel laboratorio, dovremo:

\begin{stepitemize}

\item risolvere il caso benchmark;

\item controllare fattibilit\`a, residui, slack, bound attivi e marginali;

\item variare separatamente \(\theta\), \(\phi\) e \(\psi\);

\item individuare la frontiera di fattibilit\`a rispetto a \(\theta\);

\item confrontare scenari con variazioni congiunte dei driver.

\end{stepitemize}

\vspace{0.3cm}

\textbf{Obiettivo finale:} interpretare non soltanto il valore ottimo,
ma l'intero piano ALM \(x^*,g^*,u^*\).

%TCIMACRO{\TeXButton{Transition: Box Out}{\transboxout}}%
%BeginExpansion
\transboxout%
%EndExpansion
%TCIMACRO{\TeXButton{EndFrame}{\end{frame}}}%
%BeginExpansion
\end{frame}%
%EndExpansion

\section{Caso TakeHome}

\subsection{Un nuovo problema ALM}

%******************
% Slide 19
%******************

%TCIMACRO{\TeXButton{BeginFrame}{\begin{frame}}}%
%BeginExpansion
\begin{frame}%
%EndExpansion

\QTR{frametitle}{Caso TakeHome: First Republic Bank 2023}

%TCIMACRO{\TeXButton{Transparency}{\setbeamercovered{transparent=20}}}%
%BeginExpansion
\setbeamercovered{transparent=20}%
%EndExpansion

Il caso TakeHome trasferisce lo stesso metodo di lavoro a un problema ALM
differente, ispirato al caso First Republic Bank.

%\vspace{0.4cm}

La struttura del problema presenta alcuni elementi nuovi:

\begin{stepitemize}

\item cinque classi di attivo;

\item quattro date di bilancio;

\item due possibili ondate di deposit run;

\item capacit\`a di funding limitata;

\item vincoli di ristrutturazione del bilancio.

\end{stepitemize}

%\vspace{0.4cm}

I tre driver dell'analisi sono $
\boxed{\alpha}
\qquad
\boxed{\beta}
\qquad
\boxed{\phi},
$
%\vspace{-0.1cm}
a cui corrispondono i fabbisogni:
\vspace{-0.1cm}

\[
d(\alpha,\beta)
=
(32\alpha,18,14,32\beta)'.
\]

%\vspace{0.3cm}

Due vincoli caratterizzano in particolare il caso:
\vspace{-0.1cm}
\[
x_3+x_4\ge35,
\qquad
x_5\le35.
\]

\vspace{0.2cm}

\textbf{Idea chiave:} il contesto finanziario cambia, ma la struttura
metodologica della soluzione resta la stessa.

%TCIMACRO{\TeXButton{Transition: Box Out}{\transboxout}}%
%BeginExpansion
\transboxout%
%EndExpansion
%TCIMACRO{\TeXButton{EndFrame}{\end{frame}}}%
%BeginExpansion
\end{frame}%
%EndExpansion

\subsection{Trasferire il metodo}

%******************
% Slide 20
%******************

%TCIMACRO{\TeXButton{BeginFrame}{\begin{frame}}}%
%BeginExpansion
\begin{frame}%
%EndExpansion

\QTR{frametitle}{Dal Caso Aula al Caso TakeHome}

%TCIMACRO{\TeXButton{Transparency}{\setbeamercovered{transparent=20}}}%
%BeginExpansion
\setbeamercovered{transparent=20}%
%EndExpansion

Il caso TakeHome non \`e solo una variazione parametrica del Caso Aula.

\vspace{-0.3cm}

\[
\begin{array}{c|c}
\text{Caso Aula} & \text{Caso TakeHome}\\
\hline
\text{Silicon Valley Bank} & \text{First Republic Bank}\\[0.15cm]
4\ \text{attivi} & 5\ \text{attivi}\\[0.15cm]
\theta,\phi,\psi & \alpha,\beta,\phi\\[0.15cm]
\text{stress sui fabbisogni iniziali} &
\text{due ondate di deposit run}\\[0.15cm]
\text{vincolo di concentrazione} &
\text{vincolo legacy}
\end{array}
\]

\vspace{0.1cm}

La sequenza metodologica resta invece comune:

\[
\boxed{\text{modello}}
\longrightarrow
\boxed{\text{forma solver}}
\longrightarrow
\boxed{\text{benchmark}}
\longrightarrow
\boxed{\text{controlli}}
\]

\[
\longrightarrow
\boxed{\text{sensitivit\`a e scenari}}
\longrightarrow
\boxed{\text{interpretazione}}.
\]

\vspace{0.3cm}

\textbf{Obiettivo del TakeHome:} verificare la capacit\`a di trasferire
autonomamente il metodo a una struttura finanziaria differente.

%TCIMACRO{\TeXButton{Transition: Box Out}{\transboxout}}%
%BeginExpansion
\transboxout%
%EndExpansion
%TCIMACRO{\TeXButton{EndFrame}{\end{frame}}}%
%BeginExpansion
\end{frame}%
%EndExpansion



\section{Passaggio al laboratorio}

\subsection{Sviluppo del Caso Aula in VS Code}

%******************
% Slide 21
%******************

%TCIMACRO{\TeXButton{BeginFrame}{\begin{frame}}}%
%BeginExpansion
\begin{frame}%
%EndExpansion

\QTR{frametitle}{Dal modello al laboratorio in VS Code}

%TCIMACRO{\TeXButton{Transparency}{\setbeamercovered{transparent=20}}}%
%BeginExpansion
\setbeamercovered{transparent=20}%
%EndExpansion

Abbiamo definito il problema economico, la forma matematica e la
traduzione richiesta dal solver.

\vspace{0.4cm}

Ora passiamo allo sviluppo operativo del Caso Aula in VS Code.

\vspace{0.5cm}

\[
\boxed{\text{dati}}
\longrightarrow
\boxed{\text{funzione solver}}
\longrightarrow
\boxed{\text{benchmark}}
\longrightarrow
\boxed{\text{controlli}}
\]

\[
\longrightarrow
\boxed{\text{sensitivit\`a}}
\longrightarrow
\boxed{\text{frontiera}}
\longrightarrow
\boxed{\text{scenari}}.
\]

\vspace{0.5cm}

Durante lo sviluppo, dovremo verificare continuamente la coerenza tra:

\[
\text{modello teorico}
\quad\longleftrightarrow\quad
\text{codice}
\quad\longleftrightarrow\quad
\text{output}.
\]

\vspace{0.3cm}

\textbf{Obiettivo:} costruire una soluzione computazionale
comprensibile, verificabile e interpretabile.

%TCIMACRO{\TeXButton{Transition: Box Out}{\transboxout}}%
%BeginExpansion
\transboxout%
%EndExpansion
%TCIMACRO{\TeXButton{EndFrame}{\end{frame}}}%
%BeginExpansion
\end{frame}%
%EndExpansion





\end{document}

